% Shared look for the one-page Pollis workflow handouts (same fonts and colours as the brochure).
\documentclass[10pt,a4paper]{article}

\usepackage[a4paper,left=17mm,right=17mm,top=14mm,bottom=16mm,footskip=8mm]{geometry}
\usepackage[T1]{fontenc}
\usepackage[default,semibold]{sourcesanspro}
\usepackage{microtype}
\usepackage{graphicx}
\usepackage[table]{xcolor}
\usepackage{tikz}
\usepackage[most]{tcolorbox}
\usepackage{fontawesome5}
\usepackage{enumitem}
\usepackage{fancyhdr}
\usepackage[hidelinks]{hyperref}

\graphicspath{{images/}}

% ---------------------------------------------------------------- colours (as in the brochure)
\definecolor{navy}{HTML}{1C3A5E}
\definecolor{pblue}{HTML}{3B78B5}
\definecolor{amber}{HTML}{F2B33D}
\definecolor{deepamber}{HTML}{D9902A}
\definecolor{pale}{HTML}{EEF3F9}
\definecolor{ink}{HTML}{2A2F36}
\definecolor{logonavy}{HTML}{173D63}
\definecolor{logodark}{HTML}{C98A1E}
\definecolor{logoamber}{HTML}{F2B446}
\definecolor{logocream}{HTML}{FEEFD2}
\definecolor{soft}{HTML}{5E6773}
\color{ink}

\setlength{\parindent}{0pt}
\setlength{\parskip}{3pt}
\setlist{nosep,leftmargin=*,itemsep=2pt}

\pagestyle{fancy}
\fancyhf{}
\renewcommand{\headrulewidth}{0pt}
\fancyfoot[L]{\footnotesize\color{soft}\textbf{Pollis}\enspace Bright Insights through Statistical Learning, Simulation and Optimisation}
\fancyfoot[R]{\footnotesize\color{soft}github.com/saragga/pollis/releases/latest}

% The Pollis logo, redrawn so it prints sharply at any size
\newcommand{\pollislogo}[1]{%
  \begin{tikzpicture}[x=#1/512,y=-#1/512,baseline=(current bounding box.center)]
    \fill[logonavy] (0,0) rectangle (512,512);
    \fill[logodark] (36,422) -- (104,256) -- (172,422) -- cycle;
    \fill[logoamber] (104,422) -- (203,180) -- (302,422) -- cycle;
    \fill[logocream] (203,422) -- (339,89) -- (476,422) -- cycle;
  \end{tikzpicture}}

% Top of the page: logo, name, kind of document (Workflow unless given, e.g. [Drill-Down]) and its subject;
% then the question and a one-line promise
\newcommand{\workflowhead}[4][Workflow]{%
  \thispagestyle{fancy}%
  \begin{minipage}[c]{13mm}\pollislogo{11mm}\end{minipage}%
  \begin{minipage}[c]{\dimexpr\linewidth-13mm\relax}
    {\color{navy}\Large\bfseries Pollis}\enspace{\color{soft}#1}\hfill
    {\color{deepamber}\small\bfseries\MakeUppercase{#2}}
  \end{minipage}\par\vspace{7mm}
  {\color{navy}\fontsize{22}{26}\selectfont\bfseries #3}\par\vspace{4pt}
  {\color{soft}\large #4}\par\vspace{5mm}}

% A menu path: \menu{Explore}{Access Data Libraries}... is written as \path{Explore \gt Access Data Libraries}
\newcommand{\gt}{\,{\color{deepamber}\bfseries ›}\,}
\newcommand{\menupath}[1]{{\footnotesize\bfseries\color{navy}#1}}

% A screenshot if images/<file> exists, otherwise an orange placeholder that says what to capture
\newlength{\shotheight}\setlength{\shotheight}{43mm}
\newcommand{\stepshot}[2]{%
  \IfFileExists{images/#1}{%
    \tcbox[enhanced, colback=black, colframe=black!70, boxrule=0.4pt, arc=2pt,
      left=0pt,right=0pt,top=0pt,bottom=0pt, boxsep=0pt, halign=center]{%
      \includegraphics[width=\dimexpr\linewidth-1pt\relax,height=\shotheight,keepaspectratio]{#1}}}{%
    \begin{tcolorbox}[enhanced, colback=white, colframe=orange!80, boxrule=0.6pt, arc=2pt,
      height=\shotheight, valign=center, halign=center, left=6pt, right=6pt]
      {\color{orange}\Large\faCamera}\par\smallskip
      {\color{orange}\footnotesize\bfseries Screenshot: #2}\par\smallskip
      {\color{orange!80!black}\scriptsize images/#1}
    \end{tcolorbox}}}

% One step: number, title, menu path, image file, what to capture, what it means
\newcommand{\step}[6]{%
  \begin{tcolorbox}
    {\color{deepamber}\Large\bfseries #1}\enspace{\color{navy}\large\bfseries #2}\par\vspace{1pt}
    \menupath{#3}\par\vspace{4pt}
    \stepshot{#4}{#5}\par\vspace{3pt}
    {\small #6}
  \end{tcolorbox}}

\newenvironment{steps}{%
  \begin{tcbraster}[raster columns=2, raster equal height=rows, raster column skip=4mm,
    raster row skip=4mm, enhanced, colback=pale, colframe=pale, boxrule=0pt, sharp corners,
    left=6pt,right=6pt,top=6pt,bottom=6pt, fontupper=\raggedright]}{\end{tcbraster}}

% Boxes for the bottom of the page
\newtcolorbox{pbox}[2][]{enhanced, colback=pale, colframe=amber, boxrule=0pt,
  leftrule=3pt, sharp corners, left=7pt, right=7pt, top=6pt, bottom=6pt,
  fontupper=\small\raggedright, before upper={{\color{navy}\normalsize\bfseries #2}\par\smallskip}, #1}
\newtcolorbox{nbox}[1][]{enhanced, colback=navy, colframe=navy, boxrule=0pt,
  sharp corners, left=8pt, right=8pt, top=6pt, bottom=6pt, coltext=white,
  fontupper=\small\raggedright, #1}

% A Next Steps item inside the bottom boxes
\newcommand{\nsitem}[2]{\item[\color{deepamber}\faChevronRight] \textbf{#1}: #2}
